\section{Reproducibility, code and data}
\label{app:reproducibility}

\paragraph{Code.}
All analyses are modules of one repository. Each writes one result file per subject (or per subject and repeat) in a
directory keyed by a hash of the code version and the configuration, so runs can resume and each result can be traced to
the code that produced it. The numbers in the text, tables, and figures are inserted by script from stored result files.
A completeness check lists the subjects every reported sweep must contain and stops the build of its numbers if any is
missing. The supplement holds an anonymized archive with the code, tests, and cohort cache, with no raw data,
version-control history, or author identifiers.

\paragraph{Data.}
The three cohorts are public \citep{das2025cgmacros,zhao2023shanghai,hall2018glucotypes}. Their licences are those of
the original releases, and we do not redistribute the raw files. The code reads a content-hashed cache of the parsed
cohorts. Its manifest records the source, the selection rule, and a checksum, so a machine without the raw files
reproduces the same subjects.

\paragraph{Compute.}
The analyses run on CPU. Random streams are seeded from SHA-256 hashes of the unit, so results do not depend on the number
of workers or their order. The package pins the version of every dependency, and a three-subject smoke test of the whole
pipeline runs on a laptop. Commands for every experiment are listed in the release notes of the code archive.

\paragraph{The analysis plan.}
The ten hypotheses of Table~\ref{tab:plan}, their thresholds, and the record of every later analysis, with what had been
seen when it was declared, are text files in the repository. For the most recent runs, the statement each result would
support or remove was written down before the result was read.
